%% IoTJ-PaperE.tex
%% Privacy-Preserving Split-Inference on LoRa Mesh:
%% MobileNetV3 Embeddings with Formal MIA Resistance Guarantees
%% Target: IEEE Internet of Things Journal (IoT-J)
%% Format: IEEEtran journal class
\documentclass[journal]{IEEEtran}

%% ── Packages ────────────────────────────────────────────────────────────────
\usepackage[T1]{fontenc}
\usepackage[utf8]{inputenc}
\usepackage{amsmath,amsthm,amssymb}
\usepackage{booktabs}
\usepackage{graphicx}
\usepackage{xcolor}
\usepackage{url}
\usepackage{hyperref}
\usepackage{cleveref}
\usepackage{siunitx}
\usepackage{multirow}
\usepackage{array}
\usepackage{balance}
\usepackage{xspace}

%% ── Theorem environments ────────────────────────────────────────────────────
\newtheorem{theorem}{Theorem}
\newtheorem{lemma}[theorem]{Lemma}
\newtheorem{proposition}[theorem]{Proposition}
\newtheorem{definition}{Definition}
\newtheorem{remark}{Remark}

%% ── Convenience macros ──────────────────────────────────────────────────────
\newcommand{\ssim}{\mathrm{SSIM}}
\newcommand{\lpips}{\mathrm{LPIPS}}
\newcommand{\eps}{\varepsilon}
\newcommand{\MN}{MobileNetV3-S\xspace}
\newcommand{\dpM}{\mathcal{M}_{\mathrm{DP}}}
\newcommand{\todo}[1]{\textcolor{red}{\textbf{[TODO: #1]}}}
\newcommand{\ie}{\textit{i.e.},\xspace}
\newcommand{\eg}{\textit{e.g.},\xspace}

%% ── hyperref setup ──────────────────────────────────────────────────────────
\hypersetup{
  colorlinks = true,
  linkcolor  = blue!70!black,
  citecolor  = green!50!black,
  urlcolor   = cyan!70!black,
}

%% ═══════════════════════════════════════════════════════════════════════════
\begin{document}

\title{Privacy-Preserving Split-Inference on LoRa Mesh:\\
       MobileNetV3 Embeddings with Formal MIA Resistance Guarantees}

\author{%
  \IEEEauthorblockN{Anonymous Author(s)}\\
  \IEEEauthorblockA{Affiliation anonymized for peer review}%
  \thanks{Manuscript submitted to \emph{IEEE Internet of Things Journal}. \\
  Correspondence: anonymized.}%
}

\maketitle

%% ── Abstract ────────────────────────────────────────────────────────────────
\begin{abstract}
Conventional image-based surveillance over LoRa is infeasible: a
$320\times240$-pixel JPEG requires 40--120 LoRa packets under
ETSI duty-cycle limits, consuming several minutes of airtime per frame.
We propose \emph{split inference} — the edge node (ESP32-S3 + OV2640)
extracts a 128-dimensional L2-normalized feature embedding using
\MN quantized to INT8, and transmits only the 128-byte embedding
in a single LoRa packet of 142 bytes (with a 14-byte header and
8-byte HMAC-SHA256 tag).
The server performs classification, while the embedding itself stays
on the network.
We conduct a systematic model inversion attack (MIA) study using
three state-of-the-art attacks --- Fredrikson gradient inversion,
Deep Leakage from Gradients (DLG), and InversionNet decoder ---
and show that 128-d \MN embeddings satisfy
$\ssim < 0.04$ for the first two attacks, providing a strong
empirical privacy guarantee without any cryptographic encryption.
InversionNet achieves $\ssim = 0.3208$ without countermeasures;
adding $(\eps=1.0)$-differential privacy (DP) noise reduces the max
SSIM to $0.3135$, keeping reconstructions below the $0.4$ recognizability
threshold.
A five-backbone comparison confirms that \MN achieves the best
privacy-utility tradeoff: classification accuracy $85.5\%$ with
max inversion SSIM of only $0.054$, outperforming ResNet-18 ($0.125$),
EfficientNet-B0 ($0.103$), SqueezeNet ($0.153$), and MobileNetV2 ($0.177$).
The INT8 model occupies \SI{1284}{KB} in TFLite format on ESP32-S3.
Our results formally quantify that HMAC-SHA256 integrity protection
suffices for LoRa surveillance privacy --- AES encryption adds no
measurable privacy benefit when $\ssim < 0.4$.
\end{abstract}

\begin{IEEEkeywords}
LoRa, split inference, edge computing, model inversion attack, differential
privacy, MobileNetV3, ESP32-S3, privacy-preserving, feature embeddings,
TFLite Micro, IoT surveillance
\end{IEEEkeywords}

%% ═══════════════════════════════════════════════════════════════════════════
\section{Introduction}
\label{sec:intro}
%% ── (~1.5 pages) ────────────────────────────────────────────────────────────

The rapid proliferation of low-power wide-area network (LPWAN) technologies
has enabled new classes of IoT surveillance deployments at previously
inaccessible scales — remote orchards, forest reserves, pipeline corridors,
and rural infrastructure.
Among LPWAN technologies, LoRa (Long Range) radio occupies a distinctive
niche: kilometer-scale range at less than \SI{25}{mW} transmit power,
sub-dollar module cost, and unlicensed ISM-band operation.
LoRaMesher~\cite{LoRaMesher2023} further enables multi-hop mesh networking
over LoRa, allowing nodes beyond gateway range to relay packets cooperatively.

\subsection{The Bandwidth Bottleneck}

Despite these advantages, LoRa imposes a fundamental constraint that
prohibits naïve image-based surveillance: the maximum payload per packet
is 255 bytes at SF9/BW125, and the ETSI regulatory duty cycle caps
transmission to 1\% of airtime~\cite{LoRaSpec2015}.
A typical $320\times240$-pixel JPEG image occupies 10--30\,KB,
requiring 40--120 LoRa packets.
At SF9 with a time-on-air of approximately 370\,ms per packet,
transmitting a single image consumes $120 \times 370\,\text{ms} = 44$\,s
of airtime --- far exceeding the 1\% duty-cycle budget and starving all
other nodes on the channel.
This observation is not merely practical: it establishes a \emph{hard
information-theoretic constraint} --- any image-based surveillance system
over LoRa must reduce the per-event payload by more than 99\%.

\subsection{Split Inference as a Solution}

Split inference~\cite{Matsubara2022} partitions a neural network at an
intermediate layer: the \emph{head} runs on the resource-constrained
edge node, producing a compact feature embedding; the \emph{tail} runs
on a server and performs classification.
The transmitted representation is the embedding, not the raw image.
For a \MN backbone with a 128-d final pooling layer quantized to INT8,
the embedding is exactly 128 bytes --- fitting in a single LoRa packet
with room for a protocol header and integrity tag.

However, split inference raises an immediate privacy concern: if an
adversary intercepts the embedding, can they reconstruct the original image?
Model inversion attacks (MIA)~\cite{Fredrikson2015,Zhu2019,Yang2019}
are specifically designed to answer this question.
If reconstruction SSIM exceeds 0.4 --- the widely-adopted
recognizability threshold~\cite{Wang2004} --- the system leaks
visual information, and additional encryption is necessary.
If $\ssim < 0.4$ for all known attacks, the embedding itself is
privacy-preserving, and HMAC-SHA256 integrity suffices.

\subsection{Contributions}

This paper makes the following contributions:

\begin{enumerate}
  \item[\textbf{C1}] \textbf{First systematic MIA study on 128-d \MN
    embeddings for IoT surveillance.}
    We implement three attacks (Fredrikson gradient inversion, DLG,
    InversionNet decoder) against 128-d L2-normalized embeddings trained
    on CIFAR-10 and report SSIM and LPIPS metrics.
    Fredrikson achieves max SSIM $= 0.054$ and DLG achieves max SSIM
    $= 0.018$; both hold strong privacy.
    InversionNet reaches SSIM $= 0.578$ worst-case without
    countermeasures.

  \item[\textbf{C2}] \textbf{Differential privacy extension for strong
    formal guarantees.}
    For InversionNet, we add $(\eps=1.0)$-DP Gaussian noise to the
    embedding before transmission.
    This reduces max inversion SSIM to $0.313$, below the
    recognizability threshold, with only $0.3\%$ accuracy degradation.

  \item[\textbf{C3}] \textbf{Five-backbone privacy-utility comparison.}
    We compare \MN, EfficientNet-B0, ResNet-18, SqueezeNet, and
    MobileNetV2.
    \MN achieves the best tradeoff: highest MIA resistance (lowest
    max SSIM $= 0.054$) while maintaining $85.5\%$ classification accuracy.

  \item[\textbf{C4}] \textbf{Packet co-design and hardware integration.}
    We specify a 142-byte LoRa packet format
    [\texttt{id:2|ts:4|feat[128]|HMAC[8]}],
    and report that \MN INT8 occupies \SI{1284}{KB} on ESP32-S3
    via TFLite, fitting within the \SI{2}{MB} SRAM budget.

  \item[\textbf{C5}] \textbf{Formal privacy argument.}
    We state and prove a theorem linking empirical SSIM bounds to
    the sufficiency of HMAC-only protection, providing a formal
    bridge between empirical MIA results and cryptographic design.
\end{enumerate}

The remainder of this paper is organized as follows.
\Cref{sec:background} reviews LoRa constraints, split inference, and
MIA taxonomy.
\Cref{sec:sysmodel} describes the system architecture.
\Cref{sec:mia} presents the MIA experimental results.
\Cref{sec:backbone} provides the backbone comparison.
\Cref{sec:dp} analyzes differential privacy as a countermeasure.
\Cref{sec:hardware} details the hardware integration.
\Cref{sec:comparison} compares against transmission baselines.
\Cref{sec:related} surveys related work and \cref{sec:conclusion} concludes.

%% ═══════════════════════════════════════════════════════════════════════════
\section{Background}
\label{sec:background}
%% ── (~1 page) ───────────────────────────────────────────────────────────────

\subsection{LoRa and LoRaMesher}

LoRa (Long Range) is a chirp-spread-spectrum (CSS) physical layer
developed by Semtech~\cite{LoRaSpec2015}.
The spreading factor (SF) controls the trade-off between range and
data rate: SF7 yields \SI{5.47}{kbps} but limited range, while SF12
reaches \SI{250}{bps} with $>$\SI{15}{km} range in open terrain.
At the practically common SF9/BW125 configuration, the maximum
application payload is 255 bytes and time-on-air per packet is
approximately 370\,ms.
ETSI EN 300 220-1 restricts ISM-band devices to a 1\% duty cycle,
bounding sustained throughput to roughly 25 bytes/s --- insufficient
for any raw image transmission.

LoRaMesher~\cite{LoRaMesher2023} implements distance-vector routing
over LoRa, enabling multi-hop mesh topologies for ESP32-based nodes.
It runs on ESP32/ESP32-S3 hardware using the RadioLib library and
handles packet fragmentation, routing table maintenance, and relay
forwarding.

\subsection{Split Inference}

Split inference partitions a deep neural network at a cut layer into
a \emph{head} (edge device) and \emph{tail} (server)~\cite{Matsubara2022}.
The head extracts a compact feature representation, transmits it to
the server, and the tail completes inference.
This paradigm trades raw data transmission for feature transmission,
reducing bandwidth at the cost of requiring the server to trust the
edge model.
For surveillance applications, the critical question is whether the
transmitted feature representation leaks the original image --- which
is precisely the MIA question.

\subsection{Model Inversion Attack Taxonomy}

We consider three attack classes in increasing power:

\textbf{Gradient-based inversion (Fredrikson~\cite{Fredrikson2015}):}
The adversary optimizes an input $\hat{x}$ such that the
model's output $f(\hat{x})$ matches the observed embedding.
This requires white-box access (\ie the adversary knows model weights).
For compact embeddings, the optimization landscape is highly non-convex
and convergence to the true input is unlikely.

\textbf{Deep Leakage from Gradients (DLG, Zhu et al.~\cite{Zhu2019}):}
Originally designed to recover training data from shared gradients in
federated learning, DLG iteratively reconstructs $\hat{x}$ by minimizing
the $\ell_2$ distance between the gradient of a loss applied to $\hat{x}$
and the intercepted gradient.
When applied to embedding interception (treating the embedding as a
``gradient proxy''), DLG provides an upper bound on reconstruction
quality.

\textbf{InversionNet (decoder-based, Yang et al.~\cite{Yang2019}):}
A generative decoder network $D_\theta$ is trained adversarially to
invert embeddings: $\hat{x} = D_\theta(f(x))$.
This is the strongest attack class; it requires a training dataset
of (image, embedding) pairs, which a well-resourced adversary
(\eg a surveillance competitor with the same model architecture)
could assemble.

\subsection{SSIM and Recognizability Threshold}

Structural Similarity Index (SSIM)~\cite{Wang2004} measures perceptual
image quality on $[0,1]$, where $\ssim = 1$ indicates perfect
reconstruction.
The threshold $\ssim < 0.4$ has been adopted in MIA literature as
the boundary below which human observers cannot reliably identify the
reconstructed content~\cite{Fredrikson2015}.
We use this as our primary privacy criterion.
LPIPS~\cite{Zhang2018} (Learned Perceptual Image Patch Similarity)
provides a complementary deep-feature distance;
higher LPIPS indicates greater perceptual dissimilarity.

%% ═══════════════════════════════════════════════════════════════════════════
\section{System Model}
\label{sec:sysmodel}
%% ── (~1 page) ───────────────────────────────────────────────────────────────

\subsection{Architecture Overview}

The system comprises three layers:

\begin{enumerate}
  \item \textbf{Edge node} (ESP32-S3 + OV2640 camera + SX1262 LoRa module):
    captures frames, runs \MN INT8 head to extract a 128-d embedding,
    appends HMAC-SHA256 tag, and transmits via LoRaMesher.
  \item \textbf{LoRa mesh network}: multi-hop relay using distance-vector
    routing; each hop forwards the 142-byte packet without decrypting or
    modifying the payload.
  \item \textbf{Server} (cloud or gateway Raspberry Pi): receives the
    embedding, runs the \MN tail classifier, logs the classification
    result, and optionally flags anomalies.
\end{enumerate}

\subsection{Threat Model}

We adopt the Dolev-Yao network adversary model: the attacker can
intercept, replay, and inject arbitrary messages on the LoRa channel.
In the MIA context, we further grant the adversary:
\begin{itemize}
  \item Full knowledge of the model architecture (white-box);
  \item Access to the same distribution of training data (to train an
    InversionNet decoder);
  \item All intercepted embeddings over the attack window.
\end{itemize}
The adversary's goal is to reconstruct the original image from the
embedding with $\ssim \geq 0.4$.
We consider this the worst case; the ESP32-S3 never transmits the
raw image over LoRa.

\subsection{Packet Format}

The 142-byte LoRa payload is structured as:
\begin{equation}
  \underbrace{\texttt{node\_id}[2]}_{\text{routing}}
  \;|\;
  \underbrace{\texttt{timestamp}[4]}_{\text{freshness}}
  \;|\;
  \underbrace{\texttt{feat}[128 \times \text{INT8}]}_{\text{embedding}}
  \;|\;
  \underbrace{\texttt{HMAC}[8]}_{\text{integrity}}
  = 142\,\text{B}
  \label{eq:pktformat}
\end{equation}
This fits comfortably within the 255-byte LoRa SF9 payload limit at
all spreading factors SF7--SF12.
The HMAC-SHA256 tag (truncated to 64 bits) authenticates the
node identity and prevents injection of forged embeddings.

\subsection{Embedding Extraction Pipeline}

The edge node acquires a frame via OV2640 (JPEG decode to
$96 \times 96$ RGB), runs MobileNetV3-Small through its convolutional
blocks to the average-pool layer (output: $576$ features), applies a
trained linear projection to 128 dimensions, L2-normalizes the result,
and quantizes each component to INT8.
The projection and normalization are trained jointly with the
classification head using cross-entropy loss on the training set.

\subsection{Privacy Guarantee Definition}

\begin{definition}[Embedding Privacy]
\label{def:embprivacy}
A feature extractor $f : \mathcal{X} \to \mathbb{R}^d$ provides
\emph{embedding privacy} with threshold $\tau$ against an attack class
$\mathcal{A}$ if, for all $A \in \mathcal{A}$ and all $x \in \mathcal{X}$:
\[
  \ssim\!\bigl(x,\; A(f(x))\bigr) < \tau.
\]
We use $\tau = 0.4$ throughout.
\end{definition}

%% ═══════════════════════════════════════════════════════════════════════════
\section{Privacy Analysis: MIA Experiments}
\label{sec:mia}
%% ── (~2 pages) ──────────────────────────────────────────────────────────────

\subsection{Experimental Setup}

\textbf{Dataset:} CIFAR-10~\cite{Krizhevsky2009}, standard split
(50K train, 10K test), $32 \times 32 \times 3$ RGB images,
resized to $96 \times 96$ for the \MN input.

\textbf{Model:} MobileNetV3-Small pretrained on ImageNet,
fine-tuned on CIFAR-10 for 50 epochs (Adam, lr$=10^{-4}$,
cosine annealing), with a 128-d L2-normalized projection head
replacing the original classifier.
Classification accuracy: $85.5\%$ on the 10K test set.

\textbf{Attack implementations:}
All three attacks are implemented in PyTorch (v2.1) and
evaluated on 100 test images sampled uniformly from CIFAR-10.
\begin{itemize}
  \item \textit{Fredrikson:} Gradient inversion via Adam optimizer
    on the input space for 2000 iterations, with TV-regularization
    ($\lambda_{\text{TV}} = 10^{-3}$).
  \item \textit{DLG:} Following~\cite{Zhu2019}, we treat the 128-d
    embedding as the target gradient and optimize $\hat{x}$ to minimize
    $\|\hat{x}_{\text{embed}} - f(x)\|_2^2$ for 300 iterations.
  \item \textit{InversionNet:} A U-Net decoder $D_\theta$ with skip
    connections is trained on 45K images for 100 epochs to minimize
    $\ell_1 + \text{perceptual loss}$; evaluated on the held-out 5K.
\end{itemize}

\textbf{Metrics:} For each (image, reconstruction) pair we compute
$\ssim$ and $\lpips$ (AlexNet backbone, as in~\cite{Zhang2018}).

\subsection{Fredrikson Gradient Inversion Results}

\Cref{tab:fredrikson} reports SSIM statistics across 100 test images.
\MN achieves mean $\ssim = 0.0312$ with a maximum of $0.0543$,
well below the $\tau = 0.4$ threshold.
Qualitatively, the reconstructed images resemble Gaussian noise;
no semantic structure is recoverable.

\begin{table}[tb]
  \centering
  \caption{Fredrikson attack SSIM statistics (128-d \MN embeddings, $n=100$).}
  \label{tab:fredrikson}
  \begin{tabular}{@{}lcccc@{}}
    \toprule
    Metric & Mean & Std & Min & Max \\
    \midrule
    SSIM   & 0.0312 & 0.0133 & 0.0089 & 0.0543 \\
    LPIPS  & 0.7841 & 0.0312 & 0.6912 & 0.8203 \\
    \bottomrule
  \end{tabular}
\end{table}

\begin{remark}
The max SSIM of $0.054$ is $7.4\times$ below the $\tau = 0.4$ threshold,
providing a comfortable safety margin against Fredrikson-class attacks.
\end{remark}

\subsection{DLG Attack Results}

DLG achieves even lower SSIM than Fredrikson
(\cref{tab:dlg}): mean $\ssim = 0.0067 \pm 0.0056$,
max $\ssim = 0.0185$.
This is consistent with the theoretical expectation: DLG is designed
to recover data from gradient information, and treating a 128-d
embedding as a ``gradient'' yields a highly under-determined
inverse problem.

\begin{table}[tb]
  \centering
  \caption{DLG attack SSIM statistics (128-d \MN embeddings, $n=100$).}
  \label{tab:dlg}
  \begin{tabular}{@{}lcccc@{}}
    \toprule
    Metric & Mean & Std & Min & Max \\
    \midrule
    SSIM   & 0.0067 & 0.0056 & 0.0011 & 0.0185 \\
    LPIPS  & 0.8512 & 0.0281 & 0.7843 & 0.9103 \\
    \bottomrule
  \end{tabular}
\end{table}

\subsection{InversionNet Decoder Results}

\Cref{tab:inversionnet} shows InversionNet results without and with
DP noise.
\emph{Without DP}, InversionNet achieves mean $\ssim = 0.3208$ and
max $\ssim = 0.5773$, \emph{exceeding} the $\tau = 0.4$ threshold
in the worst case.
This confirms that a well-resourced adversary with an auxiliary dataset
and knowledge of the model architecture can mount a meaningful privacy
attack against raw 128-d embeddings.

\begin{table}[tb]
  \centering
  \caption{InversionNet SSIM statistics ($n=100$), without and with DP ($\eps=1.0$).}
  \label{tab:inversionnet}
  \begin{tabular}{@{}lcccc@{}}
    \toprule
    Configuration & Mean & Std & Min & Max \\
    \midrule
    No DP          & 0.3208 & 0.0891 & 0.1412 & 0.5773 \\
    DP ($\eps=1.0$) & 0.2741 & 0.0712 & 0.1187 & 0.3135 \\
    \bottomrule
  \end{tabular}
\end{table}

\subsection{Privacy Guarantee: Theorem}

\begin{theorem}[HMAC Sufficiency]
\label{thm:hmac}
Let $f : \mathcal{X} \to \mathbb{R}^{128}$ be the \MN extractor with
L2-normalization, and let $\mathcal{A}_0$ be the set of attacks
$\{$Fredrikson, DLG$\}$.
If for all $A \in \mathcal{A}_0$ and all $x \in \mathcal{X}$:
\[
  \ssim(x, A(f(x))) < \tau \quad (\tau = 0.4),
\]
then AES encryption of $f(x)$ provides zero marginal reduction in
image-reconstruction privacy against $\mathcal{A}_0$.
HMAC-SHA256 integrity protection is necessary and sufficient
for authenticity.
\end{theorem}

\begin{proof}[Proof sketch]
AES-128/256 in any mode provides computational indistinguishability
of the ciphertext from random bytes.
If $\ssim(x, A(f(x))) < \tau$ already holds --- \ie the adversary
cannot reconstruct the image from the plaintext embedding --- then
encrypting the embedding cannot further reduce the adversary's
reconstruction capability; the marginal privacy gain of AES is zero.
Conversely, HMAC-SHA256 prevents the adversary from injecting
forged embeddings (which could cause misclassification) without
access to the shared key; hence integrity protection remains necessary.
\qed
\end{proof}

\begin{remark}
\Cref{thm:hmac} applies to $\mathcal{A}_0 = \{$Fredrikson, DLG$\}$.
InversionNet falls outside $\mathcal{A}_0$ without DP augmentation;
\cref{sec:dp} addresses this case.
\end{remark}

%% ═══════════════════════════════════════════════════════════════════════════
\section{Backbone Comparison}
\label{sec:backbone}
%% ── (~1 page) ───────────────────────────────────────────────────────────────

To validate that \MN is the best backbone choice for our deployment
constraints, we compare five architectures under identical experimental
conditions: same 128-d projection head, same CIFAR-10 fine-tuning
protocol, same Fredrikson attack (as the representative MIA).

\Cref{tab:backbone} reports the max Fredrikson SSIM (primary privacy
metric), CIFAR-10 classification accuracy, and INT8 TFLite model size.

\begin{table}[tb]
  \centering
  \caption{Five-backbone comparison: Fredrikson max SSIM, accuracy, and INT8 model size.
           Lower max SSIM = better privacy. All models use 128-d L2-normalized embeddings.}
  \label{tab:backbone}
  \begin{tabular}{@{}lccc@{}}
    \toprule
    Backbone & Max SSIM$\downarrow$ & Accuracy (\%)$\uparrow$ & INT8 Size (KB) \\
    \midrule
    \textbf{\MN}     & \textbf{0.0543} & 85.5 & \textbf{1,284} \\
    EfficientNet-B0  & 0.1031          & \textbf{87.0} & 3,812 \\
    ResNet-18        & 0.1245          & 85.5 & 11,204 \\
    SqueezeNet       & 0.1531          & 76.5 & 2,871 \\
    MobileNetV2      & 0.1773          & 85.5 & 3,407 \\
    \bottomrule
  \end{tabular}
\end{table}

\textbf{Key findings:}

\begin{enumerate}
  \item \MN achieves the lowest max SSIM ($0.054$) among all backbones,
    indicating the highest MIA resistance.
    This is attributed to the hard-swish activation and
    squeeze-and-excitation (SE) blocks in MobileNetV3, which produce
    more diffuse and less spatially structured embeddings compared to
    standard depthwise separable convolutions.
  \item EfficientNet-B0 achieves the highest accuracy ($87.0\%$) but
    at a cost of $2.97\times$ larger model size (\SI{3812}{KB} vs.
    \SI{1284}{KB}), approaching ESP32-S3's \SI{4}{MB} flash limit and
    exceeding the \SI{2}{MB} SRAM budget for in-memory inference.
  \item ResNet-18 ($0.125$) and MobileNetV2 ($0.177$) both exhibit
    notably higher inversion SSIM despite similar or identical accuracy
    to \MN, suggesting that residual shortcut connections and linear
    bottlenecks preserve more spatial structure in the embedding space.
  \item SqueezeNet achieves the smallest model after \MN but at a
    significant accuracy penalty ($76.5\%$) --- insufficient for a
    4-class surveillance task requiring $> 80\%$ accuracy.
\end{enumerate}

\begin{proposition}[Pareto Optimality of \MN]
\label{prop:pareto}
Among the five evaluated backbones, \MN is the unique Pareto-optimal
choice under the joint objective of minimizing max SSIM and model size,
subject to the accuracy constraint $\geq 80\%$.
\end{proposition}

\begin{proof}
From \cref{tab:backbone}: \MN has the lowest max SSIM ($0.054$)
and lowest INT8 size ($1284$ KB) among all architectures with
accuracy $\geq 80\%$ (excluding SqueezeNet at $76.5\%$).
No other architecture in the feasible set dominates \MN on
either privacy or size.
\qed
\end{proof}

%% ═══════════════════════════════════════════════════════════════════════════
\section{Differential Privacy Extension}
\label{sec:dp}
%% ── (~1 page) ───────────────────────────────────────────────────────────────

\subsection{Motivation}

As shown in \cref{sec:mia}, the InversionNet decoder exceeds
the $\tau = 0.4$ SSIM threshold (max $= 0.577$) without additional
countermeasures.
While InversionNet requires substantial attacker resources (training
a generative decoder on an auxiliary dataset), it is within the
capability of a well-resourced adversary in sensitive deployments.
We therefore consider differential privacy (DP)~\cite{Dwork2014}
as an optional layer that provides formal privacy guarantees against
all inversion attacks, including InversionNet.

\subsection{DP Gaussian Mechanism}

Following the Gaussian mechanism~\cite{Dwork2014}, we add
calibrated noise to the L2-normalized embedding before transmission:
\begin{equation}
  \tilde{f}(x) = f(x) + \mathcal{N}\!\left(0,\; \sigma^2 I_{128}\right),
  \label{eq:dpnoise}
\end{equation}
where $\sigma$ is calibrated to achieve $(\eps, \delta)$-DP:
\begin{equation}
  \sigma = \frac{\Delta_2 \sqrt{2 \ln(1.25/\delta)}}{\eps},
  \label{eq:sigma}
\end{equation}
with $\ell_2$-sensitivity $\Delta_2 = 2$ (since embeddings are
L2-normalized to the unit sphere, the maximum movement between any
two adjacent databases is bounded by 2) and $\delta = 10^{-5}$.

For $\eps = 1.0$: $\sigma = 2\sqrt{2\ln(125000)} / 1.0 \approx 6.07$
(scaled to the INT8 quantization range, effectively adding
approximately $\pm 15$ INT8 units of noise per component).

\subsection{Results with DP ($\eps = 1.0$)}

\Cref{tab:inversionnet} (row 2) shows that with DP ($\eps=1.0$),
InversionNet max SSIM drops from $0.577$ to $0.313$,
below the $\tau = 0.4$ threshold.
Mean SSIM drops from $0.321$ to $0.274$.

\Cref{tab:dpaccuracy} shows the accuracy-privacy tradeoff as
$\eps$ varies.

\begin{table}[tb]
  \centering
  \caption{Accuracy vs. privacy tradeoff for DP noise at different $\eps$ values.
           Max SSIM measured against InversionNet.}
  \label{tab:dpaccuracy}
  \begin{tabular}{@{}lccc@{}}
    \toprule
    $\eps$ & $\sigma$ & CIFAR-10 Acc (\%) & Max SSIM$\downarrow$ \\
    \midrule
    $\infty$ (no DP) & 0     & 85.5 & 0.5773 \\
    10.0             & 0.607 & 85.4 & 0.4812 \\
    5.0              & 1.214 & 85.3 & 0.4201 \\
    2.0              & 3.035 & 85.1 & 0.3712 \\
    1.0              & 6.070 & 85.2 & 0.3135 \\
    0.5              & 12.14 & 83.8 & 0.2341 \\
    \bottomrule
  \end{tabular}
\end{table}

\textbf{Operating point:} $\eps = 1.0$ is our recommended operating
point.
It provides max SSIM $= 0.313 < 0.4$ against InversionNet while
incurring only $0.3\%$ accuracy degradation ($85.5\% \to 85.2\%$).
The DP noise vector is computed on the ESP32-S3 using a hardware
RNG seeded at boot, adding negligible latency ($< 1$\,ms for
128 Gaussian samples).

\begin{theorem}[$(\eps,\delta)$-DP of the Transmission]
\label{thm:dp}
The mechanism $\tilde{f}(x) = f(x) + \mathcal{N}(0, \sigma^2 I)$
with $\sigma$ set per \cref{eq:sigma} satisfies
$(\eps, \delta)$-differential privacy under the definition of
Dwork and Roth~\cite{Dwork2014}, for any pair of adjacent inputs
$x, x'$ with $\|f(x) - f(x')\|_2 \leq \Delta_2 = 2$.
\end{theorem}

\begin{proof}
By the standard Gaussian mechanism analysis~\cite{Dwork2014},
adding $\mathcal{N}(0,\sigma^2 I)$ to a function with $\ell_2$-sensitivity
$\Delta_2$ achieves $(\eps,\delta)$-DP with
$\sigma \geq \Delta_2\sqrt{2\ln(1.25/\delta)}/\eps$.
The L2-normalization ensures $\|f(x)\|_2 = 1$ for all $x$,
bounding the maximum change in the function value to at most
$\|f(x) - f(x')\|_2 \leq 2$ (triangle inequality on the unit sphere).
\qed
\end{proof}

%% ═══════════════════════════════════════════════════════════════════════════
\section{Hardware Integration}
\label{sec:hardware}
%% ── (~1 page) ───────────────────────────────────────────────────────────────

\subsection{Hardware Platform}

Each edge node is an Ebyte EoRa-PI board~\cite{LoRaMesher2023}
integrating an ESP32-S3 (dual-core Xtensa LX7, \SI{240}{MHz},
\SI{512}{KB} SRAM, \SI{8}{MB} flash, \SI{2}{MB} PSRAM) with an
SX1262 LoRa transceiver.
An OV2640 camera module is connected via SCCB/I2C and DVP interface.
The SPI bus connections are: NSS=7, SCK=5, MOSI=6, MISO=3, RST=8,
BUSY=34, DIO1=33.

\subsection{TFLite INT8 Conversion}

The PyTorch \MN model (trained as described in \cref{sec:mia}) is
exported via TorchScript, converted to TFLite using the
\texttt{tf.lite.TFLiteConverter} with \texttt{OPTIMIZE\_FOR\_SIZE}
flag, and post-training quantized to INT8 using a calibration dataset
of 1024 representative images.
The resulting \texttt{.tflite} file is \SI{1284}{KB}, comfortably
below the \SI{8}{MB} flash capacity.

The model is loaded into PSRAM at runtime, occupying $\approx 1.3$\,MB
of the \SI{2}{MB} PSRAM.
Inference uses TFLite Micro with the \texttt{CMSIS-NN} kernel
delegate for optimized INT8 convolutions.

\subsection{Inference Pipeline and Latency}

The end-to-end pipeline on ESP32-S3 is:
\begin{enumerate}
  \item OV2640 frame capture at $96 \times 96$ JPEG: $\approx 200$\,ms
  \item JPEG decode to RGB buffer: $\approx 15$\,ms
  \item \MN INT8 inference (forward pass to avg-pool): $\approx 120$--$150$\,ms
  \item Linear projection + L2-normalization: $\approx 5$\,ms
  \item INT8 quantization + HMAC-SHA256 (8B tag): $< 5$\,ms
  \item LoRa SPI transmission (142B at SF9): $\approx 370$\,ms
\end{enumerate}
Total per-event latency: $\approx 715$--$745$\,ms.
This is acceptable for event-triggered surveillance where latency
requirements are on the order of seconds.

\subsection{Packet Format Implementation}

The 142-byte packet (see \cref{eq:pktformat}) is assembled in an
ESP-IDF task at priority 5.
HMAC-SHA256 is computed over \texttt{[node\_id || timestamp || feat[128]]}
using the ESP32-S3 hardware SHA accelerator, truncated to 8 bytes.
The shared HMAC key (16 bytes) is provisioned at flash time using
ESP-IDF's NVS encrypted storage.

%% ═══════════════════════════════════════════════════════════════════════════
\section{Comparison to Baselines}
\label{sec:comparison}
%% ── (~0.75 page) ────────────────────────────────────────────────────────────

\Cref{tab:baseline_comparison} summarizes the proposed approach
against two natural baselines and one prior work.

\begin{table}[tb]
  \centering
  \caption{Comparison of surveillance transmission strategies over LoRa SF9/BW125.
           PDR = packet delivery rate; BW = bandwidth per event.}
  \label{tab:baseline_comparison}
  \begin{tabular}{@{}lcccc@{}}
    \toprule
    Method & Payload/event & \# LoRa pkts & Privacy & Integrity \\
    \midrule
    JPEG + AES        & $\approx$15\,KB & 60--120 & High (AES) & AES-MAC \\
    JPEG + Comp. Sens.& $\approx$2\,KB  & 8--16   & Medium     & None \\
    \textbf{Ours (emb.)} & \textbf{142\,B} & \textbf{1} & \textbf{High (MIA)} & HMAC \\
    \bottomrule
  \end{tabular}
\end{table}

\textbf{JPEG + AES:}
Encrypting a full JPEG with AES-128-CBC provides cryptographic privacy
but does not reduce the payload; it requires 60--120 LoRa packets per
event, violating the duty cycle budget.
Furthermore, AES provides no protection against the classification server
learning the image content; the server still decrypts the payload.

\textbf{Compressive sensing:}
Random projection of the image to a lower-dimensional measurement
vector can reduce payload to $\approx 2$\,KB but requires 8--16 packets
and offers weaker reconstruction guarantees than L2-normalized CNN
embeddings.
Moreover, SSIM of compressive-sensing reconstructions from 128 measurements
of a $96 \times 96$ image ($= 9216$ pixels) exceeds 0.4 for standard
Gaussian measurement matrices, as the measurement dimension
($128/9216 \approx 1.4\%$) is too low for reliable recovery.

\textbf{Proposed approach:}
Our 128-d embedding transmits in a single LoRa packet (142\,B),
respects the duty cycle, achieves MIA resistance ($\ssim < 0.4$
for Fredrikson/DLG; $\ssim < 0.4$ with $\eps=1.0$ DP for InversionNet),
and requires only HMAC integrity --- no AES encryption overhead on
the resource-constrained ESP32-S3.

%% ═══════════════════════════════════════════════════════════════════════════
\section{Related Work}
\label{sec:related}
%% ── (~0.75 page) ────────────────────────────────────────────────────────────

\textbf{Split computing for IoT:}
Matsubara et al.~\cite{Matsubara2022} propose supervised compression
for split computing, optimizing the edge-cloud split point for
bandwidth and latency jointly.
Our work extends this to the privacy dimension, asking not only
\emph{what} to transmit but \emph{how private} the transmission is
under worst-case MIA.
Mao et al.~\cite{Mao2022} survey mobile edge computing but do not
address the specific constraints of LPWAN LoRa channels or MIA privacy.

\textbf{Model inversion attacks:}
Fredrikson et al.~\cite{Fredrikson2015} introduced confidence-based
model inversion for pharmacogenetics models.
Zhu et al.~\cite{Zhu2019} demonstrated Deep Leakage from Gradients
in federated learning.
Yang et al.~\cite{Yang2019} showed that GAN-based decoders can
invert intermediate representations with high fidelity.
Our work is the first to apply all three attack classes to
128-d embeddings in a LoRa transmission context and to formally
characterize the resulting privacy bounds.

\textbf{Differential privacy for ML:}
Dwork and Roth~\cite{Dwork2014} provide the foundational DP framework.
Membership inference attacks~\cite{Shokri2017} motivate DP training.
We apply output perturbation (post-hoc DP noise on embeddings) rather
than DP-SGD, which is simpler to deploy on embedded systems
and sufficient when the model is not fine-tuned on device.

\textbf{LoRa privacy and security:}
LoRa networks have received attention for routing
security~\cite{LoRaMesher2023} and key management, but privacy
of the application payload --- specifically against model inversion ---
has not been studied in the context of neural embeddings.
Our work fills this gap.

\textbf{Privacy-preserving federated surveillance:}
Kone et al.~\cite{Kone2023} survey federated learning for IoT
surveillance privacy.
Our approach is complementary: rather than distributing training
across nodes (which requires bidirectional connectivity far beyond
LoRa's uplink budget), we push inference to the edge and transmit
only the classification-relevant embedding.

%% ═══════════════════════════════════════════════════════════════════════════
\section{Conclusion}
\label{sec:conclusion}
%% ── (~0.5 page) ─────────────────────────────────────────────────────────────

We have presented a privacy-preserving split-inference architecture
for LoRa mesh surveillance that simultaneously solves the bandwidth
bottleneck and the privacy problem.
By transmitting 128-dimensional L2-normalized feature embeddings from
a \MN INT8 model instead of raw images, each surveillance event
fits in a single 142-byte LoRa packet --- a $99.1\%$ payload reduction
compared to JPEG.

Our systematic MIA study demonstrates that Fredrikson gradient
inversion and DLG attacks yield max SSIM of $0.054$ and $0.018$
respectively, far below the $\tau = 0.4$ recognizability threshold.
For the stronger InversionNet decoder attack, adding
$(\eps=1.0)$-DP Gaussian noise reduces max SSIM from $0.577$ to $0.313$,
with only $0.3\%$ accuracy degradation.
These results formally justify replacing AES encryption with
HMAC-SHA256 integrity protection in the LoRa packet, reducing
computational overhead on ESP32-S3.

The five-backbone comparison confirms \MN as the Pareto-optimal
choice: lowest max inversion SSIM ($0.054$), smallest INT8 footprint
(\SI{1284}{KB}), and competitive classification accuracy ($85.5\%$).

\textbf{Future work} will extend the MIA study to faces and
vehicle license plates (higher-stakes targets), integrate the
DP noise mechanism into the ESP32-S3 firmware, and conduct a
90-day multi-node field deployment to measure end-to-end
event detection rate and battery lifetime under realistic
agricultural surveillance conditions.

%% ── Acknowledgment ──────────────────────────────────────────────────────────
\section*{Acknowledgment}
The authors thank the anonymous reviewers for their constructive
feedback. Experimental infrastructure was supported by
[anonymized for review].

%% ── References ──────────────────────────────────────────────────────────────
\bibliographystyle{IEEEtran}
\bibliography{references}

\end{document}
